\section{Preliminaries and the reduction}\label{sec:preliminaries}

We use the client set \(J\), candidate set \(F\), connection cost
\(\cost\), and optimum \(\OPT_k\) defined in the introduction.
Polynomial time includes the binary encodings of all distances.
We henceforth assume \(J\ne\varnothing\) and \(1\le k\le |F|\).
This section fixes the fractional representation and states the external
reduction that will remove the final constant additive surplus.

\subsection{The fractional solution}

We use the standard relaxation
\begin{equation}\label{eq:lp}
\begin{aligned}
 \text{minimize}\quad &\sum_{j\in J}\sum_{i\in F}d(j,i)x_{ij}\\
 \text{subject to}\quad&
 \sum_{i\in F}x_{ij}=1 &&(j\in J),\\
 &0\le x_{ij}\le y_i\le1 &&(i\in F,\ j\in J),\\
 &\sum_{i\in F}y_i\le k.
\end{aligned}
\end{equation}
Its rational optimum can be found in polynomial time. For fixed openings
\(y_i\), assign each client to its nearest unit of opening mass, using a fixed
order to break ties. If \(h_j(t)\) is the distance needed to accumulate mass
\(t\), put
\[
 V_j=\int_0^1h_j(t)\,dt.
\]
Then \(L=\sum_jV_j\le\OPT_k\) for an optimal solution of~\eqref{eq:lp}.

As in LP-rounding constructions such as Charikar--Li~\cite[Section~2]{CharikarLi2012},
it is convenient to split an opening into co-located pieces so that each
client's nearest-unit assignment uses whole pieces. There are only
polynomially many needed split points: each client uses at most one partial
piece in a fixed distance order. Refine each facility simultaneously at all
the resulting cut points. Denote the whole-piece assignment of client \(j\)
by \(F_j\). For a set \(U\) of pieces, write \(y(U)=\sum_{i\in U}y_i\);
thus \(y(F_j)=1\) and \(V_j=\sum_{i\in F_j}y_id(j,i)\).

Internal pieces at one location are allowed to be distinct, so their distance
can be zero. All geometric arguments use only nonnegativity, symmetry, and
the triangle inequality. Whenever a procedure returns open pieces, merge
them by their original facility identity. Merging preserves every client's
connection cost and cannot increase the number of facilities.

\subsection{Constant-additive pseudo-approximation}

\begin{definition}
A \(c\)-additive \(\alpha\)-approximation for metric \(k\)-median returns
at most \(k+c\) facilities of cost at most \(\alpha\OPT_k\).
The additive allowance \(c\) is independent of the input instance.
\end{definition}

We use the following established result.

\begin{theorem}[Li--Svensson~\cite{LiSvensson2016}, Theorem~4 of the 2014 manuscript]
\label{thm:li-svensson}
Let \(\alpha>1\), and suppose \(\mathcal A\) is a \(c\)-additive
\(\alpha\)-approximation algorithm for metric \(k\)-median.
For every \(\eta>0\), there is an ordinary
\((\alpha+\eta)\)-approximation algorithm with running time
\(n^{O(c/\eta)}\) times the running time of \(\mathcal A\).
For fixed \(\alpha,c,\eta\), this is polynomial overhead.
\end{theorem}

The theorem permits separate client and candidate sets and applies the
pseudo-algorithm to restricted candidate instances. Our construction works
on every such instance. We will obtain a deterministic pseudo-algorithm
before applying Theorem~\ref{thm:li-svensson}; the external reduction is not
being used to convert an expected-cost statement into determinism.

For each fixed accuracy we permit all grid denominators, numbers of rounds,
and additive allowances to be constants of any size. They may depend on
that accuracy, but not on \(k\), the number of input points, their distances,
or their bit lengths. Thus it suffices to prove, for each fixed
\(\tau>0\), a deterministic \(c_\tau\)-additive
\((1+2/e+\tau)\)-approximation. Our analysis in fact bounds the
pseudo-solution against \(L\), the initial fractional cost. This does not
assert the same factor for rounding the original relaxation to exactly
\(k\) facilities: the additive allowance and the subsequent reduction
are essential.

\subsection{Probability and bookkeeping}

We use \(\E\), \(\Prob\), and \(\Var\) for expectation, probability, and variance.
Every distribution used below has finite support. Randomized procedures
first provide the estimates; Section~\ref{sec:derandomization} replaces their
choices by explicit polynomial-size lists.

Two elementary facts will be used repeatedly. First, for a nonnegative random
cost \(C\) and an event \(E\) of probability at least \(1-\rho>0\),
\begin{equation}\label{eq:conditioning}
 \E[C\mid E]\le\frac{\E C}{1-\rho}.
\end{equation}
No cost--event independence is required. Second, a finite distribution
contains an outcome whose cost is at most its expected cost. These facts
allow us to discard budget failures and select a single output by total
cost; we never require all individual client expectation bounds to hold
simultaneously at that output.

The notation \(O_\Delta(\cdot)\) means that its implicit constant may depend
on the fixed grid denominator \(\Delta\), and on no input parameter.
All unqualified constants in the geometric and analytic estimates are
absolute, unless an explicit previously fixed parameter is indicated.
